\documentclass[12pt,a4paper,oneside]{article}
\usepackage[brazil]{babel}
\usepackage[utf8]{inputenc}
\usepackage{olymp}
\usepackage{color}
\usepackage[dvipsnames]{xcolor}
\usepackage{listings}

\setcounter{problem}{4}

\begin{document}

\begin{problem}{Matriz triangular}{triangular}{2}

Em álgebra linear, uma matriz triangular superior é aquela em que todos os elementos abaixo da diagonal principal são nulos (iguais a zero). Em outras palavras, se $A$ é uma matriz triangular superior, então $a_{ij}=0, \forall i>j$.

Veja um exemplo abaixo:

$$
\begin{pmatrix}
7 & 2 & 3\\
0 & 0 & 0\\
0 & 0 & 9\\
\end{pmatrix}
$$

Faça um programa para ler uma matriz quadrada de ordem $N$ e informar se ela é ou não uma matriz triangular superior.

%\Notes
%
%Observações, se tiver.

\InputFile

Cada entrada contém apenas um caso de teste. A primeira linha contém um número $N$, indicando a ordem da matriz. Em seguida virão $N$ linhas contendo $N$ inteiros cada, indicando os elementos da matriz.  Restrição $1 \leq N \leq 1000$.

\OutputFile

Seu programa deve gerar uma única linha de saída. Se a matriz for triangular superior, escreva ``\texttt{SIM}''. Caso contrário, escreva ``\texttt{NAO}''.

%  \Notes
% Você pode escrever os resultados à medida que lê a entrada, não precisa ler a entrada toda antes.

\Example

\begin{example}
\exmp{
3
1 2 3
0 5 6
0 0 9
}{
SIM
}%
\end{example}

\begin{example}
\exmp{
4
5 0 0 0
4 5 0 0
3 4 2 0
2 9 7 5
}{
NAO
}%
\end{example}

\begin{example}
\exmp{
4
1 0 0 0
0 2 0 0
0 0 3 0
0 0 0 4
}{
SIM
}%
\end{example}

\vspace{12pt}
Dica: faça uma função booleana que recebe uma matriz e sua ordem.

\end{problem}

\end{document}
